% To allow compilation of the file
% !TeX root = ./../article.tex

\section{System Model and Attacks}

The protocol has four roles. The \emph{owner} initializes the election and the
authority registry. Each \emph{local authority} computes its local tally
off-chain and submits one encrypted tally vector. The \emph{smart contract}
acts as the public bulletin board and proof gate. A small set of
\emph{trustees} performs threshold decryption of the final aggregate.

Our assumptions are deliberately narrow. Ballots remain secret inside each
authority, local counting is outside the protocol boundary, electorate sizes
are known to the contract, and the final result must be derived from the full
set of accepted local tallies rather than from a coordinator's private view.
We add one privacy requirement that cleartext bulletin boards do not satisfy:
local tallies should remain confidential even after the election.

Within this model, we consider five threats: \emph{omission and replacement} by a
coordinator; \emph{strategic observation} before the election closes;
\emph{local retaliation} under cleartext publication; malformed encrypted
submissions that hide negative, oversized, or wrong-sum tallies; and opening
inconsistencies in the trustee registry or decryption transcript.

The protocol addresses these threats with encrypted submissions, transcript
inclusion, and a finalization rule that cancels incomplete elections: local
tallies are never revealed, only the final aggregate is opened, malformed
submissions are excluded through zero-knowledge proofs, and trustee setup and
opening are proof-checked.

The protocol has explicit limitations. A proof of \emph{well-formedness} is
not a proof that the encrypted tally is the \emph{true} local tally, so a
fully corrupted local committee remains outside scope. The artifact also uses
a dealer-generated threshold key rather than distributed key generation, and
final tally decoding after threshold opening is performed off-chain through
bounded discrete-log recovery.
